\scanpage{21}
% source lines 1--8; draft cancelled with long diagonal strokes.
\deleted{\uline{Théorème} $G$ comme dessus, $C$ un sous-grpe de Cartan, $H$ un sous-schéma en groupes de $G$, simple$/S$, à fibres connexes, \uline{contenant $C$}, $\unsure{\mathfrak h|\mathscr H}\subset\unsure{\mathfrak g|\mathscr G}$ algèbre de Lie. Alors $H$ est déterminé de façon unique par $C$ et $\mathfrak h$ et les conditions précédentes, i.e. $\forall$ autre ss-grpe $H_1$ satisfaisant aux ces conditions et tel que $\mathfrak h_1=\mathfrak h$, est égal à $H$.}

% source line 9; locally struck opening within the cancelled draft.
\deleted{Lorsque les} \deleted{Soit $T$ le tore maximal de $C$.}

% source lines 10--15
\deleted{Lorsque les caractères infinitésimaux [\unsure{pour} \unsure{sur} $\unsure{\mathfrak t|t}=\operatorname{Lie}(T)$] de $T$ opérant sur $\mathfrak g/\mathfrak h$ sont $\ne0$ sur chaque fibre, de sorte que sur chaque fibre $\mathfrak h$ est son propre normalisateur, alors par Exposé de Demazure, il s'ensuit que $H=\operatorname{Norm}_{G}(\mathfrak h)^{\unsure{0|\circ}}$, et on gagne.}

% source lines 16--21
\deleted{Dans le cas général, soit $\unsure{\mathfrak z|Z}$ le centre \unsure{réductif} de $G$, qui est $\subset T$, donc dans $H$. Ceci dit : divisant par $\mathfrak z$, on} \deleted{\unsure{donc}} \deleted{est ramené à la situation où $\mathfrak z=0$. Alors $\mathfrak t$ opère \uline{fidèlement}}
% archive sheet number: 20
